\documentclass[12pt]{article}
\usepackage[utf8]{inputenc}
\usepackage{amsmath, amssymb, amsthm}
\usepackage[spanish]{babel}
\usepackage{geometry}
\geometry{a4paper, margin=1in}

\title{Solución CP03}
\author{Carlos Javier Blanco Moreira}

\begin{document}
	
	\maketitle
	
	\section*{Introducción}
	En el presente documento se resuelven tres ejercicios correspondientes a la cuarta conferencia de Matemática Numérica. Los ejercicios seleccionados abordan la interpolación polinomial mediante distintos enfoques: el método general de resolución de sistemas, el método de Newton y el método de Lagrange. Se ha puesto especial cuidado en mostrar el desarrollo paso a paso, manteniendo un estilo claro y accesible.
	
	\section{Ejercicio 1 (Apartado a)}
	
	\subsection{Enunciado}
	Sea la función \( f \) que toma valores \( f(0) = 5 \), \( f(3) = 2 \), \( f(5) = 5 \). Halle un polinomio que interpole a la función en esos puntos, sin utilizar el método de Newton ni el método de Lagrange.
	
	\subsection{Solución}
	Los puntos dados son:
	\[
	(0,5),\quad (3,2),\quad (5,5).
	\]
	Buscamos un polinomio de grado menor o igual que 2, de la forma:
	\[
	P(x) = ax^2 + bx + c.
	\]
	Imponemos las condiciones de interpolación:
	
	\begin{itemize}
		\item \( P(0) = 5 \) implica \( c = 5 \).
		\item \( P(3) = 2 \) implica \( 9a + 3b + 5 = 2 \), es decir, \( 9a + 3b = -3 \), que simplificando da:
		\[
		3a + b = -1. \tag{1}
		\]
		\item \( P(5) = 5 \) implica \( 25a + 5b + 5 = 5 \), es decir, \( 25a + 5b = 0 \), que simplificando da:
		\[
		5a + b = 0. \tag{2}
		\]
	\end{itemize}
	
	Restando la ecuación (1) de la (2):
	\[
	(5a + b) - (3a + b) = 0 - (-1) \quad \Rightarrow \quad 2a = 1 \quad \Rightarrow \quad a = \frac{1}{2}.
	\]
	Sustituyendo en (2):
	\[
	5 \cdot \frac{1}{2} + b = 0 \quad \Rightarrow \quad \frac{5}{2} + b = 0 \quad \Rightarrow \quad b = -\frac{5}{2}.
	\]
	Por tanto, el polinomio interpolador es:
	\[
	\boxed{P(x) = \frac{1}{2}x^2 - \frac{5}{2}x + 5}.
	\]
	
	\section{Ejercicio 2}
	
	\subsection{Enunciado}
	Se sabe que la suma de los cuadrados de los \( n \) primeros números naturales es un polinomio de tercer grado en \( n \). Hállalo utilizando interpolación.
	
	\subsection{Solución}
	Definimos:
	\[
	S(n) = 1^2 + 2^2 + \cdots + n^2.
	\]
	Buscamos un polinomio cúbico:
	\[
	S(n) = an^3 + bn^2 + cn + d.
	\]
	Conocemos los valores iniciales:
	\[
	S(1) = 1,\quad S(2) = 5,\quad S(3) = 14,\quad S(4) = 30.
	\]
	Utilizaremos el método de Newton con nodos \( x_0 = 1, x_1 = 2, x_2 = 3, x_3 = 4 \).
	
	Calculamos las diferencias divididas:
	
	\[
	\begin{aligned}
		f[x_0, x_1] &= \frac{5 - 1}{2 - 1} = 4, \\
		f[x_1, x_2] &= \frac{14 - 5}{3 - 2} = 9, \\
		f[x_2, x_3] &= \frac{30 - 14}{4 - 3} = 16, \\
		f[x_0, x_1, x_2] &= \frac{9 - 4}{3 - 1} = \frac{5}{2}, \\
		f[x_1, x_2, x_3] &= \frac{16 - 9}{4 - 2} = \frac{7}{2}, \\
		f[x_0, x_1, x_2, x_3] &= \frac{\frac{7}{2} - \frac{5}{2}}{4 - 1} = \frac{1}{3}.
	\end{aligned}
	\]
	
	El polinomio de Newton es:
	\[
	S(n) = 1 + 4(n-1) + \frac{5}{2}(n-1)(n-2) + \frac{1}{3}(n-1)(n-2)(n-3).
	\]
	
	Desarrollamos paso a paso:
	
	Primero, \( 1 + 4(n-1) = 4n - 3 \).
	
	Luego, \( \frac{5}{2}(n-1)(n-2) = \frac{5}{2}(n^2 - 3n + 2) = \frac{5}{2}n^2 - \frac{15}{2}n + 5 \).
	
	Finalmente, \( \frac{1}{3}(n-1)(n-2)(n-3) = \frac{1}{3}(n^3 - 6n^2 + 11n - 6) = \frac{1}{3}n^3 - 2n^2 + \frac{11}{3}n - 2 \).
	
	Sumando todos los términos:
	\[
	\begin{aligned}
		S(n) &= \left(4n - 3\right) + \left(\frac{5}{2}n^2 - \frac{15}{2}n + 5\right) + \left(\frac{1}{3}n^3 - 2n^2 + \frac{11}{3}n - 2\right).
	\end{aligned}
	\]
	
	Agrupamos por potencias:
	\begin{itemize}
		\item \( n^3 \): \( \frac{1}{3} \).
		\item \( n^2 \): \( \frac{5}{2} - 2 = \frac{5}{2} - \frac{4}{2} = \frac{1}{2} \).
		\item \( n \): \( 4 - \frac{15}{2} + \frac{11}{3} = \frac{24}{6} - \frac{45}{6} + \frac{22}{6} = \frac{1}{6} \).
		\item Término independiente: \( -3 + 5 - 2 = 0 \).
	\end{itemize}
	
	Por tanto:
	\[
	S(n) = \frac{1}{3}n^3 + \frac{1}{2}n^2 + \frac{1}{6}n.
	\]
	
	Esta es la conocida fórmula:
	\[
	\boxed{S(n) = \frac{n(n+1)(2n+1)}{6}}.
	\]
	
	\section{Ejercicio 4}
	
	\subsection{Enunciado}
	Dada la función \( f(x) = \log_5 x \), halle un polinomio interpolador de grado 2 para la función en el intervalo \([1, 25]\). Calcule con la mayor precisión posible \( f(4) \).
	
	\subsection{Solución}
	Elegimos tres nodos en el intervalo:
	\[
	x_0 = 1,\quad x_1 = 5,\quad x_2 = 25.
	\]
	Sus imágenes son:
	\[
	y_0 = \log_5 1 = 0,\quad y_1 = \log_5 5 = 1,\quad y_2 = \log_5 25 = 2.
	\]
	Usaremos el método de Lagrange.
	
	Polinomios de Lagrange:
	\[
	\begin{aligned}
		L_0(x) &= \frac{(x-5)(x-25)}{(1-5)(1-25)} = \frac{(x-5)(x-25)}{(-4)(-24)} = \frac{(x-5)(x-25)}{96}, \\
		L_1(x) &= \frac{(x-1)(x-25)}{(5-1)(5-25)} = \frac{(x-1)(x-25)}{4 \cdot (-20)} = -\frac{(x-1)(x-25)}{80}, \\
		L_2(x) &= \frac{(x-1)(x-5)}{(25-1)(25-5)} = \frac{(x-1)(x-5)}{24 \cdot 20} = \frac{(x-1)(x-5)}{480}.
	\end{aligned}
	\]
	
	El polinomio interpolador es:
	\[
	P(x) = 0 \cdot L_0(x) + 1 \cdot L_1(x) + 2 \cdot L_2(x) = -\frac{(x-1)(x-25)}{80} + \frac{2(x-1)(x-5)}{480}.
	\]
	
	Simplificamos con denominador común \( 480 \):
	\[
	P(x) = -\frac{6(x-1)(x-25)}{480} + \frac{2(x-1)(x-5)}{480} = \frac{(x-1)}{480} \left[ -6(x-25) + 2(x-5) \right].
	\]
	
	Calculamos el corchete:
	\[
	-6(x-25) + 2(x-5) = -6x + 150 + 2x - 10 = -4x + 140 = 4(35 - x).
	\]
	
	Entonces:
	\[
	P(x) = \frac{(x-1) \cdot 4(35 - x)}{480} = \frac{(x-1)(35 - x)}{120}.
	\]
	
	Evaluamos en \( x = 4 \):
	\[
	P(4) = \frac{(4-1)(35 - 4)}{120} = \frac{3 \cdot 31}{120} = \frac{93}{120} = 0.775.
	\]
	
	El valor exacto es \( \log_5 4 = \frac{\ln 4}{\ln 5} \approx \frac{1.386294}{1.609438} \approx 0.861353 \).  
	La aproximación obtenida es:
	\[
	\boxed{P(4) \approx 0.775}.
	\]
	
\end{document}